\chapter{Presentación de resultados y análisis}
\label{cap:resultados}


Este capítulo presenta y analiza los resultados de la investigación según el diseño metodológico de la \autoref{sec:metodologia}: describe la medición de los datos y su confiabilidad, ejecuta los procesos de cálculo y validación, interpreta los hallazgos y contrasta finalmente la hipótesis. La evidencia proviene de someter a EduFEM a una batería sistemática de verificación y validación (V\&V), siguiendo la distinción metodológica clásica entre \emph{verificación} y \emph{validación} establecida en el \autoref{sec:verificacion-validacion} \autocite{oberkampf2010vv}. La exposición recorre primero la verificación del código mediante el método de soluciones manufacturadas, continúa con las pruebas de consistencia interna del modelo de datos, presenta dos validaciones independientes (la viga de Timoshenko frente a la solución analítica y a SAP2000, y la membrana de Cook como referencia histórica de bloqueo por cortante), describe la cobertura pedagógica de los módulos educativos sobre el canal de cálculo del MEF, interpreta los resultados y contrasta la hipótesis. Los datos y figuras son reproducibles mediante los scripts \texttt{tests/vv\_mms.py}, \texttt{tests/vv\_timoshenko.py} y \texttt{tests/vv\_cook.py}.

\section{Medición de los datos y validación de la medición}
\label{sec:medicion}

La variable dependiente del estudio se mide a través de tres familias de indicadores numéricos. La \emph{exactitud frente a soluciones de referencia} se cuantifica con la norma $L^2$ del desplazamiento y la seminorma $H^1$ de su gradiente, con la norma $L^2$ del campo de tensiones recuperado —el que el post-proceso muestra—, en sus formas absoluta y relativa, y con los errores relativos porcentuales en tensión y desplazamiento respecto de los valores analíticos o de referencia. La \emph{respuesta estructural} se registra como desplazamientos nodales, componentes cartesianas de tensión, tensiones principales y de von Mises, y reacciones en los apoyos. La \emph{convergencia} se mide como la tasa con que decrecen las normas de error al refinar la malla, estimada por el logaritmo en base dos del cociente de errores entre dos mallas sucesivas.

La confiabilidad de estas mediciones se asegura por tres vías. En primer lugar, las normas de error se integran con un punto de Gauss más por dirección que la rigidez ($3\times3$ para Q4, $4\times4$ para Q9) para evitar el aliasing del error con la cuadratura de la matriz y no subestimarlo artificialmente. En segundo lugar, las tensiones se evalúan en los puntos de Gauss —que en el Q4 coinciden con los puntos de Barlow de superconvergencia— y se extrapolan a los nodos con la inversa de la matriz de funciones de forma evaluada en esos mismos puntos. En tercer lugar, la corrección de ese operador de extrapolación se comprueba de forma independiente exigiendo que reproduzca exactamente los campos polinómicos de su espacio de interpolación, prueba que forma parte de la batería de regresión. Los instrumentos de medición son los guiones de verificación y validación del proyecto, deterministas y reejecutables, lo que garantiza la repetibilidad de cada dato reportado.

\section{Verificación numérica por el método de soluciones manufacturadas}

El método de soluciones manufacturadas (MMS) constituye la técnica de referencia para verificar la corrección de un código de elementos finitos, ya que permite construir una solución exacta arbitraria y derivar de ella el término fuente que la hace satisfacer la ecuación de gobierno \autocite{oberkampf2010vv,salari2000mms}. Si la implementación es correcta, el error de discretización debe disminuir con el tamaño característico de malla $h$ siguiendo las tasas asintóticas que predice la teoría de aproximación: para un elemento de orden polinómico $p$, se espera $\mathcal{O}(h^{p+1})$ en la norma $L^2$ y $\mathcal{O}(h^{p})$ en la seminorma $H^1$ \autocite{strang2008analysis,hughes2000fem,zienkiewicz2013fem}.

Se adoptó como configuración base el cuadrado unitario $[0,1]^2$ en tensión plana, con material $E = 1{,}0$, $\nu = 0{,}3$ y espesor $t = 1{,}0$. La solución manufacturada elegida fue

\begin{equation}
u_M(x,y) = \sin(\pi x)\,\sin(\pi y), \qquad v_M(x,y) = \cos(\pi x)\,\cos(\pi y),
\label{eq:mms-solucion}
\end{equation}

cuyo término fuente $\bm{f} = -\nabla\cdot\bm{\sigma}(u_M)$ se inyecta como fuerza de volumen y los valores de $u_M$ se imponen como restricciones de Dirichlet no homogéneas en el contorno. Se resolvieron mallas estructuradas de $N\times N$ elementos para $N\in\{2,4,8,16,32\}$ en cuatro configuraciones: el cuadrado unitario y un cuadrilátero distorsionado de vértices $(0;0)$, $(1;0{,}15)$, $(0{,}85;1)$ y $(0{,}1;0{,}8)$ —cuyos elementos tienen Jacobiano no constante—, cada uno en tensión plana y en deformación plana, con el término fuente deducido en cada caso de la ley constitutiva correspondiente. En todas se midieron las normas $L^2$ y $H^1$ del desplazamiento y la norma $L^2$ del campo de tensiones recuperado $\bm{\sigma}^*_h$ (extrapolado a los nodos, promediado e interpolado con las funciones de forma, tal como lo muestra el post-proceso), integradas con un punto de Gauss más por dirección que la rigidez. La \autoref{tab:mms} reporta la configuración base; las otras tres se resumen en la \autoref{tab:mms-configs} del \autoref{anexo:vyv}.

La \autoref{tab:mms} reúne los errores y las tasas de convergencia observadas para ambos tipos de elemento. La tasa se calcula como el logaritmo en base dos del cociente de errores entre dos mallas sucesivas, puesto que cada refinamiento divide $h$ a la mitad.

\begin{table}[htbp]
\centering
\caption[Errores y tasas de convergencia medidas por MMS (configuración base).]{Errores y tasas de convergencia medidas por MMS para los elementos Q4 y Q9 en la configuración base (cuadrado unitario, tensión plana): normas $L^2$ y $H^1$ del desplazamiento y norma $L^2$ del campo de tensiones recuperado $\sigma^*$.}
\label{tab:mms}
\small
\setlength{\tabcolsep}{4pt}
\begin{tabular}{c r r c r c r c}
\hline
$N$ & GDL & error $L^2$ & tasa $L^2$ & error $H^1$ & tasa $H^1$ & error $\sigma^*$ & tasa $\sigma^*$ \\
\hline
\multicolumn{8}{l}{\textbf{Q4 (bilineal, $p=1$)}} \\
2  & 18   & $2{,}559\times10^{-1}$ & ---  & $1{,}496$               & ---  & $1{,}106$              & ---  \\
4  & 50   & $7{,}241\times10^{-2}$ & 1,82 & $7{,}247\times10^{-1}$  & 1,05 & $4{,}780\times10^{-1}$ & 1,21 \\
8  & 162  & $1{,}866\times10^{-2}$ & 1,96 & $3{,}578\times10^{-1}$  & 1,02 & $1{,}686\times10^{-1}$ & 1,50 \\
16 & 578  & $4{,}699\times10^{-3}$ & 1,99 & $1{,}783\times10^{-1}$  & 1,01 & $5{,}751\times10^{-2}$ & 1,55 \\
32 & 2178 & $1{,}177\times10^{-3}$ & 2,00 & $8{,}906\times10^{-2}$  & 1,00 & $1{,}976\times10^{-2}$ & 1,54 \\
\hline
\multicolumn{8}{l}{\textbf{Q9 (bicuadrático, $p=2$)}} \\
2  & 50   & $2{,}172\times10^{-2}$ & ---  & $2{,}911\times10^{-1}$  & ---  & $2{,}282\times10^{-1}$ & ---  \\
4  & 162  & $2{,}776\times10^{-3}$ & 2,97 & $7{,}263\times10^{-2}$  & 2,00 & $6{,}076\times10^{-2}$ & 1,91 \\
8  & 578  & $3{,}479\times10^{-4}$ & 3,00 & $1{,}809\times10^{-2}$  & 2,01 & $1{,}540\times10^{-2}$ & 1,98 \\
16 & 2178 & $4{,}352\times10^{-5}$ & 3,00 & $4{,}516\times10^{-3}$  & 2,00 & $3{,}860\times10^{-3}$ & 2,00 \\
32 & 8450 & $5{,}441\times10^{-6}$ & 3,00 & $1{,}129\times10^{-3}$  & 2,00 & $9{,}654\times10^{-4}$ & 2,00 \\
\hline
\end{tabular}
\end{table}

El elemento Q4 alcanza asintóticamente tasas de $\mathcal{O}(h^{2{,}00})$ en $L^2$ y $\mathcal{O}(h^{1{,}00})$ en $H^1$, que coinciden con la predicción teórica para un elemento de orden $p=1$; el elemento Q9 alcanza $\mathcal{O}(h^{3{,}00})$ en $L^2$ y $\mathcal{O}(h^{2{,}00})$ en $H^1$, igualmente coincidentes con la teoría para $p=2$. La aproximación progresiva de las tasas a su valor asintótico --por ejemplo, el Q4 pasa de 1,82 en la malla más gruesa a 2,00 en la más fina-- es la firma esperada del régimen de convergencia. Las mismas tasas, con diferencias inferiores a 0,01, se obtienen en las otras tres configuraciones (\autoref{tab:mms-configs}): la deformación plana y la distorsión de los elementos no alteran el orden de convergencia, lo que descarta errores de orden en el ensamblaje, la integración por cuadratura de Gauss, el mapeo isoparamétrico sobre elementos no rectangulares, la matriz constitutiva de ambos estados planos y la imposición de restricciones no homogéneas.

El campo de tensiones recuperado converge a $\mathcal{O}(h^{2{,}00})$ en el Q9 y a $\mathcal{O}(h^{1{,}54})$ en el Q4. Que el Q4 supere el orden $\mathcal{O}(h^{1})$ del gradiente crudo (seminorma $H^1$) refleja la superconvergencia del muestreo en los puntos de Barlow y del promediado nodal (\autoref{sec:resolucion-tensiones}) \autocite{cook2002concepts,zienkiewicz2013fem}; que no alcance $\mathcal{O}(h^{2})$ se debe a la capa de contorno, donde los nodos reciben el promedio de un solo lado y pierden la cancelación de errores que los nodos interiores obtienen de sus cuatro elementos vecinos. Conviene señalar los límites de esta verificación: como en un campo suave los cuatro valores de Gauss de un elemento fino difieren en $\mathcal{O}(h)$, una permutación errónea de las columnas de la matriz de extrapolación apenas altera la norma $L^2$ en el régimen asintótico, aunque corrompa gravemente las tensiones nodales de una malla gruesa. Un defecto exactamente de ese tipo --la matriz $\bm{E}_{\text{Q4}}$ de la \autoref{eq:extrap-q4} escrita para otra numeración de puntos de Gauss-- existió en el motor hasta la revisión final de este trabajo, pasó inadvertido a las normas de desplazamiento y fue corregido tras detectarlo al inspeccionar la memoria de cálculo del ejemplo canónico; de ahí que la batería incorpore desde entonces la prueba de reproducción exacta de campos polinómicos, que es la que garantiza el operador de extrapolación con independencia del tamaño de malla.

La \autoref{fig:convergencia} muestra las curvas de error del desplazamiento en escala doble logarítmica, donde la pendiente de cada recta corresponde a la tasa de convergencia; la \autoref{fig:convergencia-tension} muestra la del campo de tensiones recuperado.

\begin{figure}[htbp]
\centering
\begin{minipage}{0.49\linewidth}
  \centering
  \includegraphics[width=\linewidth]{mms_convergence_l2.png}\\[2pt]
  {\small (a) Norma $L^2$ del desplazamiento.}
\end{minipage}
\hfill
\begin{minipage}{0.49\linewidth}
  \centering
  \includegraphics[width=\linewidth]{mms_convergence_h1.png}\\[2pt]
  {\small (b) Seminorma $H^1$ del desplazamiento.}
\end{minipage}
\caption[Convergencia bajo refinamiento $h$ medida por MMS.]{Convergencia bajo refinamiento $h$ medida por MMS (configuración base). Las pendientes de las rectas reproducen las tasas teóricas $\mathcal{O}(h^{p+1})$ en $L^2$ y $\mathcal{O}(h^{p})$ en $H^1$; las líneas discontinuas marcan las pendientes de referencia.}
\label{fig:convergencia}
\end{figure}

\begin{figure}[htbp]
\centering
\includegraphics[width=0.62\linewidth]{mms_convergence_stress.png}
\caption[Convergencia del campo de tensiones recuperado medida por MMS.]{Convergencia bajo refinamiento $h$ del campo de tensiones recuperado $\bm{\sigma}^*_h$ (puntos de Gauss $\to$ nodos $\to$ promedio) medida por MMS en la configuración base. Las líneas discontinuas marcan las pendientes observadas, $\mathcal{O}(h^{1{,}5})$ para Q4 y $\mathcal{O}(h^{2})$ para Q9.}
\label{fig:convergencia-tension}
\end{figure}

\section{Consistencia interna del modelo de datos}
\label{sec:consistencia-interna}

Más allá de la corrección numérica del solucionador, se verificó la robustez de las operaciones de gestión del modelo que distinguen a EduFEM como herramienta interactiva. Dos pruebas de regresión resultan especialmente relevantes.

La primera es el ciclo de transformación $\text{Q4}\rightarrow\text{Q9}\rightarrow\text{Q4}$. Al expandir un elemento bilineal a bicuadrático se generan automáticamente cinco nodos adicionales (cuatro nodos medios de arista y el centroide); al reducir de nuevo a Q4 esos nodos auxiliares deben eliminarse sin contaminar el resultado. La prueba comprueba que el campo de desplazamientos calculado sobre la malla Q4 reducida es idéntico al de la malla Q4 original dentro de la tolerancia numérica, con $\max|u_{\text{Q4,inicial}} - u_{\text{Q4,reducido}}| < 10^{-9}$. Esto confirma que las rutinas de expansión y reducción de malla no introducen deriva (\emph{drift}) numérica acumulada.

La segunda prueba aborda la separación entre identidad e índice de grado de libertad. Los identificadores de nodo son identidad pública y pueden presentar huecos tras operaciones de borrado, mientras que el índice de fila o columna en las matrices del sistema es ordinal y contiguo. La prueba resuelve un mismo problema con identificadores no contiguos $\{1,5,50,99\}$ y verifica equivalencia numérica inferior a $10^{-9}$ respecto del mismo problema con identificadores contiguos $\{1,2,3,4\}$. Este resultado valida el patrón de indexación de grados de libertad adoptado, equivalente al de los paquetes profesionales, y descarta los fallos por indexación directa que aparecían al eliminar nodos.

La tercera prueba verifica la interoperabilidad de datos que exige el sexto objetivo específico (\texttt{tests/test\_interop.py}). El ejemplo canónico exportado al paquete CSV/ZIP y reimportado en un proyecto vacío conserva sus nueve nodos, cuatro elementos, material, cargas y restricciones, y el sistema resuelto sobre el modelo reimportado reproduce los desplazamientos originales con $\max|\Delta u| < 10^{-12}$. El archivo DXF de ejemplo del proyecto —una viga de $4\times1$~m subdividida en seis cuadriláteros en la capa \texttt{FEM\_ELEMENTS}, con una capa decorativa que debe ignorarse— se importa como seis elementos y doce nodos compartidos, todos con orientación antihoraria, y su reimportación es idempotente: cero nodos y cero elementos nuevos, seis duplicados omitidos. La serialización JSON del formato nativo se verifica por ida y vuelta en \texttt{tests/test\_serialization.py}, y la generación de la memoria de cálculo en PDF, para Q4 y Q9, en \texttt{tests/test\_memoria\_calculo.py}.

\section{Validación con la viga de Timoshenko}

La validación confronta el software con una referencia externa de aceptación general. Se eligió una viga simplemente apoyada de longitud $L = 14$~m, peralte $H = 1{,}20$~m y base $b = 0{,}40$~m, sometida a una carga uniforme $q = 5000$~kgf/m, modelada en el sistema técnico (kgf, cm) con material $E = 217\,370$~kgf/cm$^2$ ($\approx 21{,}32$~GPa) y $\nu = 0{,}20$; ese módulo corresponde a la correlación módulo--resistencia $E = 15\,000\sqrt{f'_c}$ (en kgf/cm$^2$), la forma en unidades técnicas de la expresión que la norma ACI 318 propone para hormigón de peso normal, evaluada en $f'_c = 210$~kgf/cm$^2$; el análisis se realiza en régimen elástico lineal, sin considerar la fisuración del material. Conviene precisar que el valor concreto de $E$ no condiciona la validación. En elasticidad lineal con cargas prescritas y apoyos de desplazamiento nulo, escalar el módulo por un factor $\alpha$ deja el campo de tensiones inalterado y divide los desplazamientos por $\alpha$; como las tres soluciones que se comparan --EduFEM, la solución analítica y el modelo de SAP2000-- usan el mismo módulo, los errores relativos de la \autoref{tab:timoshenko-stress} y de la \autoref{tab:timoshenko-defl} son invariantes ante esa elección. El hormigón realista se adoptó únicamente para que las magnitudes del ejemplo resulten familiares al lector de ingeniería civil, no porque de él dependa el resultado. Se empleó una malla Q9 estructurada de $56\times8$ elementos (448 elementos, 1921 nodos, 3842 GDL). Las dos referencias de contraste son la solución analítica de la teoría de la elasticidad de Timoshenko y Goodier \autocite{timoshenko1970elasticity} y un modelo equivalente resuelto en SAP2000 \autocite{csi2017sap2000} con elementos de cáscara (\emph{shell}) de formulación Mindlin-Reissner. Esa formulación no coincide exactamente con el continuo bidimensional de tensión plana de EduFEM —difiere en la cinemática y en la discretización—, por lo que pequeñas discrepancias entre ambos modelos son esperables y no deben interpretarse como mayor o menor exactitud de uno sobre otro. El detalle del modelo de SAP2000, la resolución analítica de Timoshenko-Goodier y la comparación punto a punto que originan estas referencias se documentan en el \autoref{anexo:sap2000}.

La \autoref{tab:timoshenko-stress} compara la tensión normal $\sigma_x$ en tres puntos de referencia interiores de la viga.

\begin{table}[htbp]
\centering
\caption{Tensión normal $\sigma_x$ (kgf/cm$^2$) en puntos de referencia: EduFEM frente a la solución analítica de Timoshenko-Goodier y a SAP2000.}
\label{tab:timoshenko-stress}
\small
\setlength{\tabcolsep}{4pt}
\begin{tabular}{c c r r r c c}
\hline
Punto & Coord.\ (m) & EduFEM & Analítico & SAP2000 & Error anal.\ (\%) & Error SAP (\%) \\
\hline
A & $(0{,}0;\ 0{,}6)$   & $127{,}8825$ & $127{,}8542$ & $127{,}8510$ & 0,0221 & 0,0246 \\
B & $(-4{,}5;\ -0{,}3)$ & $-37{,}3410$ & $-37{,}3255$ & $-37{,}2640$ & 0,0414 & 0,2066 \\
C & $(-3{,}5;\ -0{,}2)$ & $-31{,}8077$ & $-31{,}7992$ & $-31{,}7850$ & 0,0268 & 0,0715 \\
\hline
\end{tabular}
\end{table}

Los errores de EduFEM frente a la solución analítica permanecen por debajo del $0{,}05\,\%$ en los tres puntos, con un máximo de $0{,}0414\,\%$ en el punto B. La concordancia con SAP2000 es del mismo orden, con un máximo de $0{,}2066\,\%$, también en el punto B. Ambos modelos reproducen la solución analítica dentro de márgenes muy estrechos; la pequeña diferencia entre EduFEM y SAP2000 en ese punto es atribuible a la distinta formulación (continuo de tensión plana frente a cáscara) y no permite afirmar la superioridad de una herramienta sobre la otra. La elección de los puntos A, B y C --localizados a distintas alturas y abscisas-- garantiza que la concordancia se verifique tanto en zonas de tracción como de compresión y a lo largo del peralte, donde la distribución de $\sigma_x$ varía linealmente según la teoría de vigas.

La \autoref{tab:timoshenko-defl} contrasta la deflexión vertical máxima en el centro de la luz.

\begin{table}[htbp]
\centering
\caption{Deflexión vertical máxima $|v|$ en el centro de la luz: EduFEM frente a la solución analítica con corrección de cortante.}
\label{tab:timoshenko-defl}
\begin{tabular}{l r r c}
\hline
Medición & EduFEM (cm) & Analítico (cm) & Error (\%) \\
\hline
$|v|(0,0)$ & $2{,}03460$ & $2{,}02926$ & 0,2633 \\
\hline
\end{tabular}
\end{table}

El error en deflexión es del $0{,}2633\,\%$ frente a la solución analítica que incluye la corrección de cortante de Timoshenko. Este detalle es pedagógicamente significativo: si se empleara la fórmula clásica de Euler-Bernoulli, que desprecia la deformación por cortante, el valor de referencia sería $1{,}9976$~cm y el error aparente del software ascendería al $1{,}85\,\%$. La menor discrepancia frente a la teoría de Timoshenko evidencia que el modelo de elasticidad bidimensional implementado en EduFEM captura correctamente la contribución del cortante a la flecha, contribución que es relevante en vigas de canto apreciable como la analizada, donde la relación luz-peralte es del orden de $11{,}7$. Este contraste ilustra, además, la importancia de seleccionar la referencia analítica adecuada al modelo numérico que se valida.

\autoref{fig:timoshenko-contorno} presenta el contorno de tensión normal y la configuración deformada de la viga, que permiten apreciar visualmente la distribución lineal de $\sigma_x$ en el peralte.

\begin{figure}[htbp]
\centering
\includegraphics[width=\linewidth]{fig_timoshenko_contorno}
\caption[Campo de $\sigma_x$ sobre la viga de Timoshenko deformada ($56\times8$ Q9).]{Post-proceso de EduFEM: campo de tensión normal $\sigma_x$ sobre la configuración deformada de la viga de Timoshenko ($56\times8$ Q9). La gradación lineal de $\sigma_x$ en el peralte --tracción (rojo) en la fibra inferior, compresión (azul) en la superior y eje neutro (verde)-- coincide con la predicción de la teoría de vigas.}
\label{fig:timoshenko-contorno}
\end{figure}

\section{Validación con la membrana de Cook y el bloqueo por cortante}

La membrana de Cook es un caso de prueba clásico de la literatura de elementos finitos, que debe su nombre a R. D. Cook y resulta especialmente sensible al bloqueo por cortante (\emph{shear-locking}) por combinar flexión, cortante y distorsión geométrica en un cuadrilátero trapezoidal \autocite{cook2002concepts}. El dominio es un trapecio de vértices $(0,0)$, $(48,44)$, $(48,60)$ y $(0,44)$, empotrado en el borde izquierdo y cargado con una fuerza tangencial total unitaria en el borde derecho; el material es adimensional con $E = 1{,}0$, $\nu = 1/3$ y $t = 1{,}0$ en tensión plana. La cantidad de interés es el desplazamiento vertical $u_y$ en el punto $(48,52)$, para la que se adopta el valor de referencia $u_y^{\text{ref}} = 23{,}96$, de uso convencional para este caso.\footnote{A diferencia de la viga de Timoshenko, este caso no tiene solución analítica cerrada: su valor de referencia es un límite de convergencia y no un dato exacto, por lo que conviene tratarlo como tal y no como una constante tomada de una fuente. En lugar de descansar en el valor publicado, este trabajo lo corrobora con evidencia propia: la secuencia Q9 ($23{,}925 \to 23{,}949 \to 23{,}961$) todavía se mueve en la malla más fina, y su extrapolación de Richardson \autocite{oberkampf2010vv} apunta a $\approx 23{,}97$, consistente con el valor adoptado; el orden observado, próximo a uno, refleja la singularidad de tensiones en las esquinas reentrantes del empotramiento. Se adopta $23{,}96$ como referencia con la salvedad de que su incertidumbre, del mismo orden de magnitud que la diferencia entre las dos mallas más finas, es comparable al error de las mallas Q9 más finas: por debajo del $0{,}1\,\%$ las diferencias no son significativas.}

La \autoref{tab:cook} reúne la convergencia bajo refinamiento $N$ para ambos elementos.

\begin{table}[htbp]
\centering
\caption[Convergencia del desplazamiento $u_y$ en $(48,52)$ para la membrana de Cook.]{Convergencia del desplazamiento $u_y$ en $(48,52)$ para la membrana de Cook. Referencia adoptada $u_y^{\text{ref}} = 23{,}96$, con una incertidumbre del orden de 0,05 por ciento.}
\label{tab:cook}
\begin{tabular}{c r r r r r r}
\hline
 & \multicolumn{3}{c}{\textbf{Q4}} & \multicolumn{3}{c}{\textbf{Q9}} \\
$N$ & GDL & $u_y$ & error (\%) & GDL & $u_y$ & error (\%) \\
\hline
2  & 18   & $11{,}845$ & $-50{,}563$ & 50   & $23{,}289$ & $-2{,}802$ \\
4  & 50   & $18{,}299$ & $-23{,}626$ & 162  & $23{,}840$ & $-0{,}502$ \\
8  & 162  & $22{,}079$ & $-7{,}850$  & 578  & $23{,}925$ & $-0{,}144$ \\
16 & 578  & $23{,}430$ & $-2{,}210$  & 2178 & $23{,}949$ & $-0{,}044$ \\
32 & 2178 & $23{,}818$ & $-0{,}594$  & 8450 & $23{,}961$ & $+0{,}003$ \\
\hline
\end{tabular}
\end{table}

El contraste entre ambos elementos es ilustrativo. El elemento Q4 exhibe el fenómeno de bloqueo por cortante con claridad: en la malla más gruesa subestima la respuesta en más de un $50\,\%$, y aun con $N=32$ (2178 GDL) conserva un error del $-0{,}594\,\%$. Su convergencia es lenta y monótona por debajo de la referencia, comportamiento característico de un elemento que resulta excesivamente rígido cuando debe representar flexión mediante una interpolación bilineal incapaz de reproducir el modo de curvatura pura \autocite{cook2002concepts,bathe2014fem}. El elemento Q9, en cambio, converge con rapidez y prácticamente libre de bloqueo: con $N=8$ (578 GDL) ya alcanza un error del $-0{,}144\,\%$, y con $N=16$ (2178 GDL) lo reduce a $-0{,}044\,\%$, dentro de la incertidumbre de la propia referencia.

La comparación a igualdad de grados de libertad es el argumento pedagógico central. Con los mismos 2178 GDL, el Q9 ($N=16$) logra un error de $-0{,}044\,\%$ mientras que el Q4 ($N=32$) se queda en $-0{,}594\,\%$: la diferencia de precisión a igualdad de grados de libertad es de un orden de magnitud (trece veces) a favor del elemento de mayor orden, y no depende de la incertidumbre de la referencia, que solo afecta a la tercera cifra del error del Q9. Esta evidencia empírica permite al estudiante constatar de primera mano la ventaja del elemento de mayor orden frente al bloqueo, un fenómeno que de otro modo permanece abstracto.

\autoref{fig:cook-convergencia} muestra las curvas de convergencia de $u_y$ frente al número de grados de libertad para ambos elementos.

\begin{figure}[htbp]
\centering
\includegraphics[width=0.72\linewidth]{cook_convergence.png}
\caption{Convergencia de la membrana de Cook. El Q9 alcanza la referencia con muchos menos grados de libertad que el Q4, que sufre bloqueo por cortante.}
\label{fig:cook-convergencia}
\end{figure}

La \autoref{tab:resumen-q4q9} sintetiza la comparación global entre ambos elementos a lo largo de las tres baterías de prueba, consolidando la evidencia que sustenta la decisión de ofrecer ambos tipos en la herramienta.

\begin{table}[htbp]
\centering
\caption{Síntesis comparativa Q4 frente a Q9 en las baterías de verificación y validación.}
\label{tab:resumen-q4q9}
\begin{tabular}{l c c}
\hline
Aspecto & Q4 (bilineal) & Q9 (bicuadrático) \\
\hline
Tasa $L^2$ (MMS)              & $\mathcal{O}(h^{2{,}00})$ & $\mathcal{O}(h^{3{,}00})$ \\
Tasa $H^1$ (MMS)              & $\mathcal{O}(h^{1{,}00})$ & $\mathcal{O}(h^{2{,}00})$ \\
Cook $N=8$ (162 / 578 GDL)    & $22{,}079$ ($-7{,}850\,\%$) & $23{,}925$ ($-0{,}144\,\%$) \\
Cook $N=16$ (578 / 2178 GDL)  & $23{,}430$ ($-2{,}210\,\%$) & $23{,}949$ ($-0{,}044\,\%$) \\
Timoshenko $\sigma_x$ (error máx.) & --- & 0,0414\,\% \\
Comportamiento bajo distorsión & bloqueo parcial & sin bloqueo apreciable \\
\hline
\end{tabular}
\end{table}

\section{Resultados del software y cobertura del canal de cálculo del MEF}

Más allá de la corrección numérica, EduFEM se propuso cubrir de extremo a extremo el procedimiento del MEF mediante ocho módulos educativos (M0 a M7) que materializan, en orden pedagógico, cada eslabón de la cadena de cálculo. Cada módulo es una \emph{capa superpuesta} interactiva que ilumina sobre la malla real el punto donde ocurre el cálculo, evitando la abstracción de una ventana desconectada del modelo \autocite{lee2015interactive,bishay2020teaching}.

La cobertura del canal de cálculo por los módulos es completa hasta el ensamblaje; la recuperación de tensiones, último eslabón del canal, no tiene un módulo propio y se expone en el post-proceso mediante la sonda puntual —con sus modos crudo y suavizado—, la vista tridimensional y la memoria de cálculo, que desarrolla la extrapolación y el promediado con sustitución numérica. El módulo M0 atiende la calidad de malla en el pre-proceso, evaluando el Jacobiano escalado y la compacidad (\emph{stretch}) con los umbrales bibliográficos de Verdict \autocite{stimpson2007verdict}, mientras que los módulos M1 a M7 recorren la fase de proceso en el orden canónico del método —del mapeo isoparamétrico al ensamblaje global, según el mapeo módulo--etapa de la \autoref{tab:modulos}—. Ese orden no es casual: refleja la dependencia lógica entre conceptos, pues la cadena cinemática $\bm{N}\rightarrow\bm{J}\rightarrow\bm{B}$ precede a la constitutiva, y la deformación $\bm{\varepsilon} = \bm{B}\,\bm{u}$ antecede a la tensión $\bm{\sigma} = \bm{D}\,\bm{\varepsilon}$, en consonancia con la presentación clásica de la literatura \autocite{bathe2014fem,cook2002concepts}.

Cada módulo que exhibe formulación matemática incorpora un conmutador Fórmula~$\leftrightarrow$~Valores que bifurca entre la narrativa simbólica en LaTeX y la evaluación numérica en el punto seleccionado. Este mecanismo concreta una secuencia de aprendizaje deliberada: el estudiante lee primero \emph{qué} se calcula y después confronta el número real, replicable incluso en una hoja de cálculo. La selección directa en el lienzo, en lugar de controles redundantes, mantiene una única fuente de verdad para el elemento bajo análisis: el panel no puede quedar mostrando un elemento distinto del que está resaltado en la malla.

La fase de post-proceso ofrece la inspección del campo de resultados mediante contornos de tensión, configuración deformada, una vista tridimensional del campo y un sondeo puntual que despliega el círculo de Mohr del estado tensional en cualquier punto del dominio. Como complemento impreso, la Memoria de Cálculo genera un documento que reproduce el paso a paso simbólico y numérico del modelo concreto del usuario, desde las funciones de forma evaluadas en los puntos de Gauss hasta la recuperación de tensiones y su extrapolación nodal. Así, los módulos enseñan los conceptos en abstracto y la Memoria valida el cálculo sobre el caso particular, cerrando el ciclo entre teoría y aplicación.

\section{Interpretación de los resultados}

\subsection{Alcances}

La triple batería de V\&V respalda la fiabilidad del núcleo numérico. La verificación por MMS muestra que, en tensión y deformación plana y sobre mallas regulares y distorsionadas, el desplazamiento converge con las tasas teóricas y el campo de tensiones recuperado lo hace a $\mathcal{O}(h^{1{,}5})$ en Q4 y $\mathcal{O}(h^{2})$ en Q9, lo que descarta errores de orden en el ensamblaje, la integración, la imposición de restricciones y la cadena de recuperación de tensiones. Las dos validaciones independientes contrastan el software con referencias de naturaleza distinta: una solución analítica cerrada en el caso de Timoshenko, con un error del $0{,}04\,\%$ en la tensión normal $\sigma_x$, del $0{,}26\,\%$ en la flecha y de hasta el $2{,}9\,\%$ en la tensión cortante, y una referencia histórica sensible al bloqueo en el caso de Cook, donde la herramienta reproduce tanto la patología del Q4 como la convergencia del Q9. Frente a SAP2000 las diferencias son del $0{,}21\,\%$ en $\sigma_x$ y de hasta el $0{,}56\,\%$ en desplazamientos; en la componente cortante el modelo de cáscara queda más cerca de la solución analítica ($0{,}16\,\%$) que EduFEM. Para los problemas de referencia estudiados, el error en las magnitudes primarias --tensión normal y flecha-- se mantiene por debajo del $0{,}3\,\%$ frente a la solución analítica y del $0{,}6\,\%$ frente al modelo de SAP2000; las componentes secundarias de tensión, recuperadas con un orden de convergencia menor y de magnitud mucho más pequeña, requieren mallas más finas para alcanzar el mismo nivel relativo.

\subsection{Comparación cualitativa con herramientas existentes}

Frente a los paquetes comerciales como SAP2000, EduFEM no compite en cobertura de tipologías estructurales, escalabilidad ni prestaciones de productividad; su valor reside en la transparencia del cálculo. Mientras que un paquete profesional opera como una caja negra optimizada para la producción, EduFEM expone cada matriz intermedia, cada punto de Gauss y cada operación de ensamblaje al estudiante. En relación con las herramientas educativas reportadas en la literatura --simuladores de armaduras, visualizadores de los modos propios de la matriz de rigidez o entornos de procesamiento interactivo de ecuaciones del MEF \autocite{bishay2020teaching,lee2015interactive,lee2015eigenmodes}--, la contribución distintiva de EduFEM es la cobertura integral del canal de cálculo de elementos continuos bidimensionales con elementos isoparamétricos Q4 y Q9, articulada con una generación automática de memoria de cálculo que reproduce, sobre el modelo del propio usuario, el procedimiento con sustitución numérica paso a paso —un recurso que el usuario puede rehacer manualmente y contrastar, y que no se documenta entre los antecedentes educativos revisados (\autoref{sec:antecedentes})—. Esta orientación responde directamente a las necesidades de formación en MEF identificadas en un estudio con expertos de la industria y de la academia, que reclama de los egresados no solo la destreza operativa sino la capacidad de planificar, verificar y validar sus análisis \autocite{perezsantiago2023fem}. La elección de un entorno íntegro en Python sobre las bibliotecas científicas de uso extendido \autocite{harris2020numpy,virtanen2020scipy} favorece, además, la inspección y extensión del código por parte de docentes y estudiantes.

\subsection{Limitaciones observadas}

El alcance del software está acotado de manera deliberada a la elasticidad lineal plana con elementos cuadriláteros Q4 y Q9, según se delimita en la Introducción. Esas restricciones, coherentes con el objetivo pedagógico, privilegian la claridad conceptual sobre la generalidad, pero limitan la aplicabilidad a casos de ingeniería más allá de los didácticos.

Desde el punto de vista numérico, el elemento Q4 arrastra el bloqueo por cortante documentado en la \autoref{tab:cook}, que aunque se aprovecha como recurso pedagógico, exige al usuario refinar la malla o recurrir al Q9 para obtener resultados precisos en problemas dominados por flexión. No se implementaron técnicas de mitigación como la integración reducida selectiva o la formulación $\bar{\bm{B}}$, que se reservan como línea de extensión futura. En la formulación de deformación plana, además, la matriz constitutiva degenera hacia la incompresibilidad cuando el coeficiente de Poisson se aproxima a $0{,}5$, régimen en el que aparece el bloqueo volumétrico; el software lo señala visualmente pero no lo resuelve.

En cuanto a la escalabilidad, conviene precisar el estado real del motor numérico. La matriz de rigidez se almacena en formato disperso --se ensambla por tripletes de coordenadas (COO) y se convierte a filas comprimidas (CSR)-- y el sistema reducido se resuelve mediante una factorización LU dispersa directa, tal como se detalló en el \autoref{cap:marco_teorico}. Esta arquitectura está implementada, no meramente proyectada, y permite abordar mallas muy superiores a las de tamaño didáctico empleadas en las validaciones. La \autoref{tab:tiempos} cuantifica el costo de ensamblaje y solución para la membrana de Cook Q9 a refinamientos crecientes, junto con la memoria que ocuparía la matriz de rigidez en formato denso frente al disperso.

\begin{table}[htbp]
\centering
\caption[Tiempos de ensamblaje y solución, y memoria de la matriz de rigidez (Cook Q9).]{Tiempo de ensamblaje y de solución del sistema disperso, y memoria de la matriz de rigidez en formato denso frente a disperso (CSR), para la membrana de Cook Q9 a refinamientos crecientes. Medido en un equipo de escritorio con procesador Intel (microarquitectura Ivy Bridge, circa 2012) bajo Windows 10 y Python 3.13, con el motor vectorizado por lotes en NumPy puro que incorpora el ejecutable distribuido; no requiere ninguna dependencia adicional. Reproducible con \texttt{tests/bench\_timing.py}.}
\label{tab:tiempos}
\begin{tabular}{r r r r r}
\hline
GDL & $t_{\text{ensamblaje}}$ (s) & $t_{\text{solución}}$ (s) & $\bm{K}$ densa (MB) & $\bm{K}$ dispersa (MB) \\
\hline
578   & 0,004 & 0,007 & 3    & 0,21  \\
2178  & 0,013 & 0,026 & 38   & 0,81  \\
4802  & 0,028 & 0,061 & 184  & 1,81  \\
8450  & 0,049 & 0,127 & 571  & 3,20  \\
18818 & 0,111 & 0,329 & 2833 & 7,19  \\
33282 & 0,191 & 0,721 & 8862 & 12,77 \\
\hline
\end{tabular}
\end{table}

Las mediciones evidencian que la malla más exigente de las validaciones (8450~GDL) se ensambla y resuelve en menos de dos décimas de segundo, y que el motor procesa problemas de más de 33\,000~GDL en menos de un segundo, todo ello sobre hardware de más de una década de antigüedad. El contraste de memoria es igualmente ilustrativo: a 33\,282~GDL la matriz densa requeriría cerca de 8{,}9~GB, mientras que su representación dispersa ocupa menos de 13~MB --un ahorro de casi tres órdenes de magnitud que es justamente lo que hace viable el cálculo en equipos modestos--. La limitación de escalabilidad, por tanto, no reside en el almacenamiento ni en el rango de tamaños propio del uso educativo bidimensional, sino que aparece solo en el extremo de mallas masivas: el tiempo de solución crece más rápido que el de ensamblaje a medida que el llenado (\emph{fill-in}) de la factorización directa se vuelve dominante. El ensamblaje ya se realiza vectorizado por lotes y la factorización aplica un ordenamiento de mínimo grado sobre $\bm{A}^{\top}+\bm{A}$ que reduce el llenado: frente al ordenamiento genérico COLAMD de SuperLU, medido con el mismo guion sobre el sistema reducido, el tiempo de factorización y solución se reduce 1,7 veces con 2178 GDL, 2,0 veces con 8450 GDL y 2,9 veces con 33\,282 GDL, de modo que las oportunidades de optimización que persisten para problemas de gran porte —factorizaciones especializadas para sistemas definidos positivos y solucionadores iterativos precondicionados— se plantean como línea de trabajo futuro y se detallan entre las recomendaciones. Finalmente, la validación presentada se apoya en tres casos cuidadosamente seleccionados por su valor canónico; una batería más amplia de casos de referencia reforzaría la confianza en el comportamiento del software ante geometrías y condiciones de carga más diversas.

\section{Validación de la hipótesis}
\label{sec:validacion-hipotesis}

Los resultados desarrollados permiten contrastar la hipótesis de diseño y responder las preguntas de investigación planteadas en la Introducción y sistematizadas en la matriz de consistencia (\autoref{tab:consistencia}).

La \textbf{hipótesis de diseño} —la factibilidad de construir un software que exponga el canal de cálculo completo del MEF en elasticidad plana de forma transparente, interactiva, numéricamente verificable y coherente con el flujo de trabajo profesional— queda confirmada al verificarse, sobre el artefacto construido, cada uno de esos atributos: el motor recorre la cadena completa de la formulación isoparamétrica, la interfaz la organiza en las tres fases canónicas sobre un único lienzo, los módulos educativos la exponen paso a paso y la batería de V\&V acredita su verificabilidad numérica.

La pregunta \textbf{PI-1}, sobre la precisión verificable del motor, se responde afirmativamente: la verificación por soluciones manufacturadas reprodujo las tasas de convergencia teóricas en los dos estados planos y sobre mallas distorsionadas, y la validación con la viga de Timoshenko arrojó, frente a la solución analítica, un error del $0{,}04\,\%$ en la tensión normal $\sigma_x$ y del $0{,}26\,\%$ en la flecha, y diferencias de hasta el $0{,}56\,\%$ frente a SAP2000. La pregunta \textbf{PI-3}, sobre la evidencia de fenómenos numéricos, también se responde de forma positiva: la membrana de Cook reprodujo con claridad el bloqueo por cortante del elemento Q4 frente a la convergencia del Q9, cuantificado a igualdad de grados de libertad. La pregunta \textbf{PI-2}, sobre la organización de la interfaz y de los módulos, se atiende por diseño y queda documentada en el \autoref{cap:diseno} y en la cobertura del canal de cálculo descrita en este capítulo.

Con ello queda respondido el problema científico planteado en la Introducción: un software educativo que construye el canal de cálculo con las mismas funciones que exponen los módulos y la memoria, lo organiza en las tres fases canónicas sobre un lienzo único y somete su motor a verificación y validación con criterios fijados a priori hace observable y contrastable cada etapa del procedimiento, que es el efecto que el problema pedía. El objeto de estudio —el análisis de medios continuos en elasticidad plana por el MEF— quedó caracterizado numéricamente en las secciones precedentes, y el campo de acción —su enseñanza— recibe el aporte en forma de los atributos verificados del artefacto.

La única afirmación que no se somete a contrastación empírica es el impacto educativo de la herramienta sobre el aprendizaje real de los estudiantes, cuya medición —mediante un estudio con grupos y pruebas de conocimiento— excede el alcance temporal de este trabajo y se formula como recomendación de continuación.
